%% v4/sec/platform.tex — §4 Platform: Safe-Bridge (TODAES long-form)
%%
%% Provenance: verbatim mechanical carry of paper/sec/platform.tex
%% (2026-05-16) under Phase-1 task 91240545. v3 source is frozen.
%%
%% PHASE-2 expansion targets (from outline_v4.md §4 plan, ~2.5-3.0pp):
%%   - 4.1 Threat model — expand actor catalog (curious-but-passive
%%     baseline + active-MITM out-of-scope justification)
%%   - 4.2 Architecture (Fig. 1) — enlarged caption, per-crossing
%%     numbered walk-through
%%   - 4.3 Four scrub layers — each gets motivation paragraph +
%%     alternative considered + chosen-mechanism rationale
%%   - 4.4 Five hard invariants — current "summary" para becomes
%%     dedicated subsection with per-invariant code anchor
%%   - 4.5 Three-axis scrub-audit methodology — canonical-test
%%     walkthrough with assertion-by-assertion explanation

\section{Platform: Safe-Bridge}
\label{sec:platform}

\subsection{Trust boundary and threat actor}
\label{sec:platform-threat}

The benchmark runs across eleven cloud LLM endpoints drawn from
seven provider families, with six US-based and five China-based
checkpoints (full per-provider list at
\S\ref{sec:method-protocol}). Following
the v2 threat-model framing, we treat every such provider as
\textbf{curious-but-passive}: it sees every prompt and every
tool-result string we hand it, may log or train on that traffic,
and can craft tool-call arguments designed to exfiltrate foundry
content via its own sandbox---but is not assumed to mount active
man-in-the-middle attacks against the Cadence host. Active MITM,
OS-level insider threat, and direct PDK theft are out of scope.
The protected asset is foundry NDA-bound PDK content in three
forms: BSIM4 / BSIM-CMG model parameters, process-stack layer
data, and proprietary cell topology and naming.
%% PHASE-2: expand actor catalog + active-MITM out-of-scope justification.

Beyond the trust-boundary mechanics, of the eleven checkpoints,
the five OpenAI-compatible vendors
(Moonshot Kimi K2.5, MiniMax M2.7, OpenAI GPT-5.5 / GPT-5.4-mini,
Xiaomi MiMo v2.5-pro, DeepSeek V4-pro / V4-flash) surface
per-iteration \texttt{reasoning\_tokens} via
\texttt{completion\_tokens\_details}; the Anthropic and Gemini
families do not in v3 (see \S\ref{sec:method-threats} Threat~4).

\subsection{Architecture --- Fig.~\ref{fig:arch}}
\label{sec:platform-arch}

Figure~\ref{fig:arch} shows the Safe-Bridge topology: a Python
agent process calls a \filename{safe\_bridge.py} wrapper that talks
to Cadence Virtuoso IC23.1 over the \texttt{virtuoso-bridge-lite}
RAMIC TCP channel on \texttt{localhost:65061}. Six numbered
crossings mark the trust boundary---tool call, allow-listed SKILL
expression, TCP payload, scoped query, PDK reference, redacted
return.
%% PHASE-2: per-crossing numbered walk-through; enlarged caption.

\begin{figure}[h]
  \centering
  %% NB: figs/safe_bridge_arch.svg must be converted to PDF/PNG
  %% before pdflatex build. See repo: paper/figs/safe_bridge_arch.svg.
  %% v4 graphicspath resolves via {../figs/}{./figs/}; bare basename.
  \includegraphics[width=0.6\linewidth]{safe_bridge_arch}
  \caption{Safe-Bridge topology and six numbered trust-boundary
    crossings. The verifiable NDA-safe substrate that enables the
    \S\ref{sec:methodology} benchmark.}
  \label{fig:arch}
\end{figure}

\subsection{Four scrub layers}
\label{sec:platform-scrub}

Safe-Bridge inserts four defense layers, each labelled on
Fig.~\ref{fig:arch}:
%% PHASE-2: each layer gets motivation + alternative considered + chosen mechanism rationale.

\begin{enumerate}
  \item \textbf{Python parameter-name whitelist + identifier /
    numeric-literal regex.}
    \filename{\_is\_allowed\_param\_name}
    (\filename{safe\_bridge.py:2310-2356}) admits only the
    per-project YAML allow-list (defaulted to
    \texttt{[w, l, nf, m, multi, wf, r, c, nfin, fingers, idc, vdc]}),
    passing \filename{\_SAFE\_PARAM\_NAME\_RE}
    (\filename{safe\_bridge.py:186}) and never
    \filename{\_FOUNDRY\_LEAK\_RE}. A complementary numeric-atom
    regex (\filename{safe\_bridge.py:35-37}) blocks SKILL injection
    at the value boundary.

  \item \textbf{20-name SKILL entrypoint allow-list (Cadence side).}
    \filename{\_ALLOWED\_SKILL\_ENTRYPOINTS}
    (\filename{safe\_bridge.py:125-176}); arbitrary
    \texttt{load() / hiOpenLib / dbOpenCellViewByType} rejected at
    \filename{\_check\_skill\_entrypoint}
    (\filename{safe\_bridge.py:2149-2212}). Nested-call grep
    extends the guarantee to embedded subexpressions; only the
    pure data-constructor \texttt{list} is whitelisted as a
    nested call (\filename{\_ALLOWED\_SKILL\_NESTED},
    \filename{safe\_bridge.py:181}).

  \item \textbf{Return-path scrubber.}
    \filename{\_scrub} (\filename{safe\_bridge.py:255-277})
    applies \filename{\_FOUNDRY\_LEAK\_RE}'s foundry-token deny-list
    and absolute-path regexes to every string field of every SKILL
    return value (15 anchored return sites) before it crosses back to
    the agent loop.

  \item \textbf{Reasoning-trace scrub.}
    \texttt{reasoning\_content} / \texttt{thinking} is scrubbed at
    every LLM-client sink (per-vendor reply sites enumerated in
    \S\ref{sec:method-audit}) via the shared \filename{scrub} wrapper
    (\filename{safe\_bridge.py:280-290}) before the trace re-enters
    conversation history.
\end{enumerate}

\subsection{Five hard invariants (summary)}
\label{sec:platform-invariants}

Five composable invariants---PDK-token redaction, parameter-name
whitelist, SKILL-entrypoint allow-list (with nested-call
extension), inline-upload directory containment, and op-point key
whitelist (\filename{\_SAFE\_OP\_POINT\_KEYS},
\filename{safe\_bridge.py:41-56})---collectively close the
trust-boundary tiers summarised in Table~\ref{tab:tiers}, so a
misconfigured \texttt{allowed\_params} cannot reopen a closed
channel; per-invariant code anchors are in the supplementary
appendix.
%% PHASE-2: promote each invariant to its own paragraph with code anchor + composition argument.

\subsection{Trust-boundary status}
\label{sec:platform-status}

\begin{table}[t]
  \centering
  \caption{Trust-boundary tier status under the
    curious-but-passive threat model. Tiers~1--3 (replay) are
    enforced by the four scrub layers
    (\S\ref{sec:platform-scrub}) and five composable invariants
    (\S\ref{sec:platform-invariants}); first-emission and
    side-channel coverage are out of scope.}
  \label{tab:tiers}
  \footnotesize
  \begin{tabular}{c l l}
    \toprule
    \textbf{Tier} & \textbf{Attack class} & \textbf{Status} \\
    \midrule
    1 & Naive prompt injection            & enforced \\
    2 & Schema / SKILL injection          & enforced \\
    3 & Reasoning-trace re-entry (replay) & enforced$^{\dagger}$ \\
    3 & Reasoning-trace first-emission    & out of scope \\
    4 & Timing / error-class side channel & partial \\
    \bottomrule
  \end{tabular}
  \\[0.3em]
  \begin{minipage}{0.95\linewidth}\footnotesize
  $^{\dagger}$Replay closed via scrub layer~4,
  \S\ref{sec:platform-scrub}.
  \end{minipage}
\end{table}

The first-emission row remains out of scope under the
curious-but-passive model---provider-side request logs capture
unscrubbed reasoning bytes on first emission, but we do not claim
coverage of near-active providers. Closing it requires a hash-only
thinking channel or fully on-prem inference
(\S\ref{sec:conclusion} future work).
%% PHASE-2: 4.5 Three-axis scrub-audit methodology — new subsection with canonical test walkthrough.
